%% Section 1: Introduction

\section{Introduction}
\label{sec:intro}

In late October 2022, Nigeria's central bank announced the withdrawal and
redesign of its currency notes, together with strict daily caps on cash
withdrawals. By January 2023, a rolling shortage of the replacement notes had
produced hour-long queues outside commercial banks, paralysed point-of-sale
terminals, and darkened markets where almost every transaction had always
settled in cash. Prices for essential goods spiked and some small enterprises closed
their doors. In February 2023 the Supreme Court ordered the old notes to remain
legal tender, and the central bank progressively released more new currency
over the following months. The episode, the \emph{naira crunch}, lasted
approximately six months before cash supply normalised.

The crunch was an abrupt, policy-induced cash supply shock where the stock of
physical currency available for transactions fell sharply while the demand for
cash-settled transactions did not. Its severity, however, varied substantially
across Nigeria's thirty-six states and the Federal Capital Territory
(\textsc{fct}), for a reason not immediately obvious at
the time as some states had higher pre-shock adoption of digital
payment technology. In those states, enterprises and households could substitute
to mobile money wallets, USSD bank transfers, and fintech applications when
ATMs ran dry. In states where most transactions still settled exclusively in
cash, limited substitute existed.

This paper exploits that cross-sectional heterogeneity to estimate the real
costs of a cash supply shock and the extent to which digital financial
infrastructure acts as a buffer. The identification relies on three features
of the episode. First, the \CBN{} announcement on October~26, 2022 was
genuinely unanticipated, making the pre-shock distribution of fintech
adoption a predetermined variable uncontaminated by anticipation of the shock.
Second, that distribution is measured using the Nigeria Living Standards
Phone Survey (\NLPS) Round~6 which was conducted on October~15, 2022, eleven days
before the announcement, providing a state-level index of digital payment
penetration constructed from four digital-payment service categories.\footnote{%
  The index uses four digital-payment categories: mobile money, commercial
  bank accounts, mobile and \textsc{ussd} banking, and mobile banking
  applications. Four further categories (pension/insurance, securities,
  microfinance, and credit cards) are excluded because their penetration is
  below 4\% in all states and they do not constitute a meaningful cash
  substitute during the crunch; see Appendix~Table~\ref{atab:alt_treatment}
  for robustness to alternative index definitions.
  Section~\ref{sec:data} details the construction.}
Third, the parallel-trends assumption is validated using monthly composites of
nighttime-light radiance from the \textsc{nasa} \VIIRS{} Black Marble satellite
product for all 774 Nigerian local government areas (January~2019 to
September~2022), providing 45 consecutive pre-shock months over which
high- and low-fintech states can be shown to have followed statistically
identical trajectories prior to the crunch. Section~\ref{sec:design} details
the identification strategy and pre-trend validation.

The primary evidence comes from firm-level data.and the main findings are as follows.
Enterprises in states one standard deviation above the pre-shock fintech
mean were 1.7 percentage points more likely to remain in operation at the
February 2023 peak crunch and 7.7 percentage points more likely
to remain in operation fourteen months later in April 2024.
Employment was also significantly insulated as high-fintech states had
0.29 more enterprise workers during the peak crunch. 
Effects are concentrated in rural firms and small enterprises below-median
size, the units most dependent on daily cash flow.

These results contribute to three literatures. First, they extend
\citet{chodorow2020cash} to the mobile-money studies. \citet{chodorow2020cash} shows that
Indian districts with fewer bank branches lost more economic activity during
the 2016 demonetisation; the relevant margin in Sub-Saharan Africa is not bank
branches, but digital payment technology.
This paper identifies a \emph{digital payment substitution channel} that the
India paper's design, which relies on bank-deposit variation, cannot isolate.
The magnitude of the enterprise-survival estimates implies that closing the
fintech adoption gap between the 10th- and 90th-percentile Nigerian states
would have reduced firm exit at the peak crunch by approximately 5 percentage
points, a substantial share of the 16-percentage-point decline in enterprise
activity observed between Round~5 and Round~7.

Second, this paper provides evidence for a digital-payment substitution mechanism 
emphasized in the mobile-money literature. By enabling transactions and transfers to be 
executed without physical currency, digital payment technologies reduce reliance on 
cash and may attenuate the effects of cash-supply disruptions \citep{jack2014risk}. 
The present study quantifies the extent of this insulation 
in LMICs using plausibly exogenous pre-determined variation in fintech adoption. 


Third, the analysis contributes to the emerging methodological literature on
identifying monetary transmission in data-sparse economies by leveraging alternative proxies 
for economic activity 
\citep[see][]{henderson2012measuring,donaldson2016view}. High-frequency satellite
imagery combined with granular phone survey data can recover real monetary
transmission at subnational resolution where the formal national accounts
provide limited quarterly or monthly benchmarks. The satellite--survey combination
is directly replicable for other African central bank episodes.

Section~\ref{sec:context} describes the naira crunch and Nigeria's payment
landscape. Section~\ref{sec:literature} surveys the related literature.
Section~\ref{sec:theory} presents the theoretical framework.
Section~\ref{sec:data} introduces the data. Section~\ref{sec:design}
presents the research design and validates the identifying assumption using
the 45-month \NTL{} pre-trend window. Section~\ref{sec:results_micro} reports
the firm-level results. Section~\ref{sec:longrun} examines longer-run effects
using the 2025 World Bank Enterprise Survey and the April 2024 \NLPS{} round.
Section~\ref{sec:conclusion} concludes with policy implications.
